\documentclass[11pt,a4paper]{article}
\usepackage[utf8]{inputenc}
\usepackage[T1]{fontenc}
\usepackage[spanish,es-tabla]{babel}
\usepackage{geometry}
\usepackage{graphicx}
\usepackage{float}
\usepackage{hyperref}
\usepackage{tabularx}
\usepackage{booktabs}
\usepackage{array}
\geometry{margin=2.5cm}
\hypersetup{colorlinks=true,linkcolor=blue,urlcolor=blue}
\setlength{\parindent}{0pt}
\setlength{\parskip}{0.55em}
\newcolumntype{Y}{>{\raggedright\arraybackslash}X}
\newcommand{\code}[1]{\texttt{\detokenize{#1}}}
\newcommand{\evidencia}[2]{%
\begin{figure}[H]\centering
\IfFileExists{#1}{\includegraphics[width=0.88\linewidth]{#1}}{%
\fbox{\begin{minipage}[c][3cm][c]{0.82\linewidth}\centering Pendiente de captura: \texttt{\detokenize{#1}}\end{minipage}}}
\caption{#2}
\end{figure}}

\begin{document}
\begin{titlepage}
\centering
\vspace*{1cm}
{\Large Pontificia Universidad Javeriana}\\
{\large Fundamentos de Ingeniería de Software}\\[1.5cm]
{\Huge \textbf{DriveControl / AutoMinder Enterprise}}\\[0.4cm]
{\Large Equipo SYNTIX TECH}\\[1.2cm]
{\LARGE \textbf{Milestone 2 - Gestión básica de flota}}\\[1.2cm]
\begin{tabular}{ll}
Profesor: & Luis Gabriel Moreno Sandoval\\
Fecha: & 19 de mayo de 2026\\
Repositorio: & \url{https://github.com/puj-course/FIS_2610_3517_G4}\\
\end{tabular}
\vfill
\end{titlepage}

\tableofcontents
\newpage

\section{Resumen ejecutivo del milestone}
El Milestone 2 consolidó la gestión básica de flota con interfaces de conductores y vehículos, validación RUNT, patrones de diseño y pruebas de flujos críticos. GitHub live registra 37 issues cerradas, 14 HU cerradas estrictas y 1 PR asociado.

\section{Objetivos del milestone}
\begin{itemize}
\item Implementar flujos de conductores y vehículos.
\item Integrar validación RUNT simulada.
\item Aplicar separación de responsabilidades y patrones.
\item Fortalecer pruebas funcionales de flota.
\end{itemize}

\section{Alcance funcional}
El alcance cubre gestión básica de flota, validaciones, componentes visuales y soporte de pruebas. No incluye cierre CI/CD, SMS ni dashboard final.

\section{Evidencia ágil}
\begin{tabularx}{\linewidth}{l r Y}
\toprule
Elemento & Valor & Evidencia\\
\midrule
HU cerradas & 14 & Issues con patrón estricto de HU.\\
Issues cerradas & 37 & Milestone cerrado en GitHub.\\
PRs & 1 & PR \#233 mergeado.\\
Commits & 164 & Ventana S5-S8 en Git local.\\
\bottomrule
\end{tabularx}

\evidencia{../../Entrega-Final/evidencias/01_agile_milestones_semestre.png}{Milestones del semestre.}
\evidencia{../../Entrega-Final/evidencias/02_agile_project_board.png}{Tablero GitHub Project.}

\section{Métricas del milestone}
\begin{tabularx}{\linewidth}{Y r}
\toprule
Métrica & Valor\\
\midrule
HU cerradas & 14\\
Issues cerradas & 37\\
Issues abiertas & 0\\
Commits por ventana de sprint & 164\\
PRs asociados & 1\\
PRs mergeadas & 1\\
\bottomrule
\end{tabularx}

Participación por issues asignadas: \code{solonlosada2006} 18, \code{Sarm-m} 9, \code{Juserora} 6, \code{juanvargax} 5 y \code{samuelfl680} 3.

\section{Evaluación del producto}
Se produjo gestión básica de flota, validación RUNT y pruebas de flujos críticos. El producto avanzó en mantenibilidad por separación de responsabilidades y reducción de código mezclado.

\section{Evaluación del proceso}
Mejoró la división en sub-issues, pero faltó mapa de dependencias e interoperabilidad entre entornos. GitHub sostuvo la trazabilidad del avance.

\section{Evaluación del esfuerzo}
El esfuerzo fue sostenido con 164 commits en S5-S8. La participación fue visible por issues asignadas, con carga fuerte en DevOps/flota y coordinación Scrum.

\section{Postmortem del milestone}
\subsection{Qué salió bien}
Se cerraron las issues, se fortaleció la arquitectura y se validaron flujos críticos.

\subsection{Qué salió mal}
Hubo problemas por nombres, alias, diferencias de entorno y componentes construidos antes de infraestructura base.

\subsection{Causas raíz}
La falta de convenciones de nombres y de entorno homogéneo causó bugs de importación e integración.

\subsection{Impacto}
El impacto fue retrabajo, menor visibilidad de integración y validaciones tardías.

\subsection{Acciones correctivas}
Estandarizar nombres, mapear dependencias y validar flujos en más de un entorno.

\subsection{Acciones preventivas y responsables sugeridos}
Configuration Manager debe definir convenciones; DevOps debe guiar entorno; QA debe mapear dependencias; Scrum Master debe cuidar sub-issues.

\section{Retrospectiva estrella de mar}
\subsection{Comenzar a hacer}
\begin{itemize}
\item Estandarizar convenciones de nombres mediante reglas compartidas.
\item Definir rutas y menús antes de construir componentes dependientes.
\item Recolectar evidencia técnica durante el sprint.
\end{itemize}

\subsection{Más de}
\begin{itemize}
\item Comunicar cambios en esquemas y contratos.
\item Revisar dependencias técnicas en kick-off.
\item Realizar commits frecuentes para reducir conflictos.
\end{itemize}

\subsection{Seguir haciendo}
\begin{itemize}
\item Dividir el trabajo en sub-issues atómicas.
\item Aplicar patrones de diseño cuando reduzcan deuda real.
\item Usar pruebas de flujos críticos.
\end{itemize}

\subsection{Menos de}
\begin{itemize}
\item Asumir compatibilidad entre entornos sin verificación.
\item Posponer feedback visual hasta el cierre.
\item Depender de datos manuales no reproducibles.
\end{itemize}

\subsection{Dejar de hacer}
\begin{itemize}
\item Construir componentes sobre infraestructura inexistente.
\item Duplicar archivos sin revisar componentes existentes.
\item Ignorar conflictos de alias como problemas individuales.
\end{itemize}

\section{Plan de mejora del siguiente ciclo}
\begin{tabularx}{\linewidth}{Y Y Y Y}
\toprule
Acción & Responsable & Evidencia esperada & Criterio de éxito\\
\midrule
Definir convenciones & Configuration Manager & Regla o documento & Menos errores de importación\\
Mapear dependencias & QA Lead & Checklist & Menos retrabajo\\
Validar entornos & DevOps Engineer & Guía reproducible & Resultado consistente\\
Mantener sub-issues & Scrum Master & Issues enlazadas & Trazabilidad clara\\
\bottomrule
\end{tabularx}

\section{Conclusión del milestone}
M2 fue finalizado y fortaleció la flota. Su aprendizaje central fue que arquitectura y proceso necesitan convenciones y entornos homogéneos.

\section{Anexos o evidencias pendientes}
\begin{itemize}
\item Pendiente de captura: flujos RUNT por escenario.
\item Pendiente de auditoría: relación exacta commit por HU.
\end{itemize}
\end{document}
